\section{Proofs}\label{app:proofs}

\textbf{Proposition~\ref{prop:min}.} Equation~\eqref{eq:decomp} is the chain rule $\KL(Q\Vert P)=\KL(Q_Z\Vert P_Z)+\E_{Z\sim Q}\KL(Q_{X|Z}\Vert P_{X|Z})$ for the coarse-graining $X\mapsto Z=\mathbf{1}[E]$. The conditional terms are non-negative and vanish for $Q(\cdot\mid E)=P(\cdot\mid E)$ and $Q(\cdot\mid E^c)=P(\cdot\mid E^c)$. With $Q(E)=\epsilon$ this leaves $\KL(Q\Vert P)\ge d_{\mathrm{Ber}}(\epsilon\Vert p)$ with equality exactly for $Q^\star(x)=P(x)\epsilon/p$ on $E$ and $P(x)(1-\epsilon)/(1-p)$ off $E$.

\textbf{Proposition~\ref{prop:causal}.} The ratio of $Q^\star/P$ on $E$ to that off $E$ is $\epsilon(1-p)/[p(1-\epsilon)]=2^{-\lambda}$, so $Q^\star(x)\propto P(x)2^{-\lambda\mathbf{1}[x\in E]}$. Marginalizing sequentially, the transition of $Q^\star$ at a history $h$ is proportional to $P(v\mid h)\,\E_P[2^{-\lambda\mathbf{1}[E]}\mid h,v]=P(v\mid h)[1-(1-2^{-\lambda})H(h,v)]$.

\textbf{Proposition~\ref{prop:accounting}.} Let $V(h)=\sum_vp(v\mid h)H(h,v)$ for a history that has not looped, with $V=0$ at the horizon, and $\Phi_t=\mathbf{1}[\text{looped by }t]+\mathbf{1}[\text{not looped by }t]V(h_t)$. Then $\Phi_0=P(E)$, $\Phi_N=\mathbf{1}[E]$, and under a Markov controller $\E_Q[\Phi_{t+1}-\Phi_t\mid h_t]=\sum_vq(v\mid h_t)H(h_t,v)-V(h_t)=-\Delta(h_t)$ for $h_t$ not yet looped. Summing over $t$ and taking expectations gives (a). For (b), $\Phi_t$ is a martingale under $P$ with $\Phi_0=p$ and $\Phi_N=\mathbf{1}[E]$, so $p(1-p)=\mathrm{Var}(\mathbf{1}[E])=\sum_t\E_P[\mathrm{Var}(\Phi_{t+1}\mid h_t)]$, and for a history that has not looped $\mathrm{Var}(\Phi_{t+1}\mid h_t)=\sigma^2(h_t)$.

\textbf{Proposition~\ref{prop:waste}.} Tilting $p$ by $e^{-\lambda H}$ gives $\Delta=\lambda\sigma^2+O(\lambda^2)$ and $\KL=\frac{\lambda^2\sigma^2}{2\ln2}+O(\lambda^3)$ bits, hence $\KL=\Delta^2/(2\sigma^2\ln2)\,(1+o(1))$ as $\Delta/\sigma\to0$ (iii). For (i), Cauchy--Schwarz gives $D^2=\big(\sum_h(\sqrt{w_P}\Delta/\sigma)(\sqrt{w_P}\sigma)\big)^2\le\big(\sum_hw_P\Delta^2/\sigma^2\big)\big(\sum_hw_P\sigma^2\big)=2\ln2\,K_2\,p(1-p)$, using Proposition~\ref{prop:accounting}(b); equality holds iff $\Delta/\sigma\propto\sigma$, i.e.\ $\lambda_h$ constant, and then $K_2=D^2/(2\ln2\,p(1-p))$, which is the second-order expansion of $d_{\mathrm{Ber}}(p-D\Vert p)$. For (ii), with $u_h=\Delta/\sigma^2=\lambda_h$ and $\pi_h=w_P\sigma^2/(p(1-p))$ one has $D=p(1-p)\,\E_\pi[u]$ and $\sum w_P\Delta^2/\sigma^2=p(1-p)\,\E_\pi[u^2]$, so $K_2/K^\star=\E_\pi[u^2]/\E_\pi[u]^2=1+\mathrm{CV}^2$. Restricting the Cauchy--Schwarz inequality to the touched set $A$ gives $\E_\pi[u]^2\le\varphi\,\E_\pi[u^2]$, hence $K_2/K^\star\ge1/\varphi$. The weights $w_P$ are used in place of $w_Q$; they differ by a relative $O(D)$, which does not change the leading order.
